\section*{Part VI: Construction and Capital-Project Execution}
\addcontentsline{toc}{section}{Part VI: Construction and Capital-Project Execution}


\subsection{43. The life cycle of a capital project}

Concept $\rightarrow$ feasibility $\rightarrow$ land/use diligence $\rightarrow$ brief $\rightarrow$ schematic design $\rightarrow$ detailed design $\rightarrow$ permits $\rightarrow$ procurement $\rightarrow$ mobilization $\rightarrow$ construction $\rightarrow$ testing and commissioning $\rightarrow$ occupancy approvals $\rightarrow$ operational handover $\rightarrow$ defects period $\rightarrow$ long-term maintenance. These are linked but not interchangeable. Starting physical construction before critical design and permits may convert a manageable delay into expensive rework.

Create a gate document and named approval authority at each stage. No construction contract should be treated as a creative substitute for a resolved project brief.


\subsection{44. Project delivery models}

\textbf{Design-bid-build} separates consultant design and contractor execution; offers design control but may introduce interface risk. \textbf{Design-build} puts more design and construction responsibility under 1 contractor, but client requirements must be precise. \textbf{EPC/turnkey} may be appropriate for complex works with defined outputs, subject to contract and market practice. \textbf{Construction management} offers flexibility but requires strong owner capability.

No delivery model eliminates owner risk. The best form depends on design maturity, complexity, contractor market, capital availability and the institute's management competence.


\subsection{45. Who belongs on the owner's side?}

At material scale appoint a project sponsor, owner's project manager, architect, structural and MEP engineers, quantity surveyor/cost consultant, legal and permitting advisers, safety oversight and independent quality checks. Segregate cost certification and payment approval where possible. The contractor's project manager is not an independent fiduciary of the owner.

A founder with strong software and AI expertise should not assume competence in building contracts, soil mechanics or construction claims. Buy independent expertise to reduce expensive errors.


\subsection{46. Cost estimation has maturity levels}

At concept stage, use broad ranges and explicit assumptions. At schematic design, estimate by work packages and quantities. Before tender, require a detailed bill of quantities or equivalent cost plan with independently checked rates. After contractor selection, maintain an approved baseline and change log. If the scope is uncertain, the contingency must reflect uncertainty; a fixed percentage copied from another project is not risk management.

Distinguish \textbf{construction cost}, \textbf{project capex}, \textbf{financing cost}, and \textbf{opening operating cash}. They are not interchangeable.


\subsection{47. A hypothetical \textnormal{₹}25-crore 1st facility}

[Model - phase-1 academic/research facility, not a site quote]

\begin{table}[t]
\centering
\small
\caption{Hypothetical \textnormal{₹}25-Crore Phase-1 Facility Capital Budget}
\label{tab:ch34_facility_budget}
\begin{tabularx}{\columnwidth}{X r}
\toprule
\textbf{Use Category} & \textbf{\textnormal{₹} crore} \\
\midrule
Site / legal / acquisition & 5.00 \\
Site preparation / off-site connections & 1.50 \\
Main buildings / civil and MEP & 10.00 \\
Lab fit-out and initial equipment & 2.50 \\
Power, water, network, safety & 1.50 \\
Architects, engineers, permits, supervision & 0.90 \\
Construction contingency & 2.10 \\
Pre-opening and commissioning & 0.50 \\
1st operating reserve, separately identifiable & 1.00 \\
\midrule
\textbf{Sources required} & \textbf{25.00} \\
\bottomrule
\end{tabularx}
\end{table}
The model contains \textbf{\textnormal{₹}24 crore of asset/project-related uses and \textnormal{₹}1 crore opening reserve}. A real institution may require far more operating liquidity; calculate it from payroll, enrollment and research commitments. Never treat contingency as permission for undocumented change orders.


\subsection{48. Sources and uses must match over time}

A \textnormal{₹}25-crore project with \textnormal{₹}25 crore of letters of interest is not funded. Separate cash on hand, binding equity/sponsor commitments, executed loan facilities, restricted grants, conditional pledges, vendor credit and in-kind support. Then map each against the monthly spending schedule. An equipment donation cannot pay a contractor's civil invoice, and a donor's 3-year pledge may arrive after a year-1 payment is due.

Maintain a month-by-month funding gap schedule and refuse to award contracts before the critical path is financeable.


\subsection{49. Procurement and tendering}

Publish a clear scope, qualifying criteria, drawings, performance standards, measurement rules, payment milestones, retention terms, quality testing, delay remedies, warranties, insurance, health and safety provisions, subcontracting limits, claims processes and change-control regime. Compare bids on \textbf{evaluated total value}, not only lowest price. An abnormally low bid may imply under-scoping, inferior quality or future variation claims.

For material contracts, require competitive bids or a documented sole-source justification and independent approval. Record conflicts where the builder or land seller is related to an officer or investor.


\subsection{50. Change orders can defeat a sound budget}

Changes may arise from design evolution, hidden conditions, regulatory comments, user requests and contractor errors. Classify each variation by cause, schedule impact, cost, contingency draw and approving authority. Require a written instruction and priced assessment before work where practicable. Unauthorized ``small changes'' can accumulate to a major financial exposure.

The founder should be able to answer weekly: contract sum, approved variations, pending claims, forecast cost at completion, contingency remaining, and financing gap.


\subsection{51. Construction cash-flow forecasting}

A construction budget is not a cash forecast. Track mobilization advances, measured progress payments, retention, certification lag, deposits for equipment, taxes, lender drawdowns, donor receipts, interest during construction, and defects/warranty holdback. Tie draw requests to independently verified work and lender documentation.

\textbf{[Model]} If the project incurs \textnormal{₹}3 crore in a quarter but the lender reimburses eligible certified work after 60 days, the institution may need bridging liquidity well beyond its nominal equity share. A thin funding plan can force delayed contractor payments and thereby make the original schedule unattainable.


\subsection{52. Schedule risk and critical path}

Permits, land access, long-lead electrical equipment, utility connections, monsoon, contractor procurement, specialized lab fit-out and commissioning can control completion. A Gantt chart is useful only if it identifies dependencies and critical paths. Ask what activities can proceed without irreversible exposure when 1 permit is delayed.

Maintain a downside schedule and formal escalation triggers. A 1-year delay can cost rent, salaries, finance interest and lost enrollment-not merely 12 months of inconvenience.


\subsection{53. Commissioning and practical completion}

A facility is not ready because the building appears complete. Require testing of electrical, water, HVAC, fire detection and suppression, lifts, IT, lab-specific equipment, safety systems, accessibility and operational procedures. Verify occupancy/completion permissions as required. Train the facilities team, receive as-built drawings, warranties, O\&M manuals, spare parts, and defect lists.

Do not admit learners or operate hazardous laboratories before the relevant legal approvals and safe operational conditions exist.
